% GENERATED FILE. Edit canonical lessons, then run npm run generate:books.
% canonical-source-hash: fnv1a64:2c13013c369ed8af
% canonical-lessons: RU-C470-styd, RU-C470-styditsya, RU-C470-smushchatsya, RU-C470-krasnet, RU-C470-zastenchivyy

\chapter{Shame and Embarrassment}
\label{ch:ru-shame-and-embarrassment}
\hlchaptermodality{470}

\begin{chapteropening}
\textbf{By the end of this chapter:} \emph{I can talk about shame and embarrassment.}

Complete the last lesson of chapter 470: I can talk about shame and embarrassment.
\end{chapteropening}

\section[\texorpdfstring{\ru{стыд} (styd)}{стыд (styd)}]{\ru{стыд} (styd) — shame}

\label{lesson:RU-C470-styd}



\begin{quote}
Before the new one: say the Russian for envy, then the Russian for to envy.
\end{quote}

\subsection*{You'll want to know: \ru{стыд}}
\textbf{\ru{стыд}} (\emph{styd}) — \textquotedblleft{}shame\textquotedblright{}.

\begin{quote}
\ru{От} \ru{стыда} \ru{он} \ru{не} \ru{мог} \ru{сказать} \ru{ни} \ru{слова}.
\end{quote}

\begin{quote}
\emph{Ot stydá on ne mog skazát ni slóva.}
\end{quote}

\begin{quote}
He couldn't say a word for shame.
\end{quote}

Masculine; genitive \ru{стыда}. \emph{\ru{Мне} \ru{стыдно}} — I feel ashamed.

\begin{cousinweb}
\emph{\ru{стыд}} is inherited from Proto-Slavic \emph{styd\ru{ъ}} (EtymDB); it is usually linked to the root of \emph{\ru{стыть}} 'to grow cold'.
\end{cousinweb}

\begin{grammarlens}[title={What you've built}]
One more word for telling, explaining and describing at B1.
\end{grammarlens}

\subsection*{Your turn}
\begin{itemize}
\raggedright
  \item \emph{Say it:} \emph{styd}
  \item \emph{Say it:} \emph{styd}, once more
  \item \emph{Say it:} a sentence of your own with \emph{styd}
  \item \emph{Recall:} say the Russian for envy, then the Russian for to envy, then say \emph{styd} again
\end{itemize}

\begin{tcolorbox}[breakable,colback=teal!4,colframe=teal!35!black,title={Before you move on}]
What does \ru{стыд} mean? (\textquotedblleft{}Shame\textquotedblright{}.) Say it once more.
\end{tcolorbox}
\section[\texorpdfstring{\ru{стыдиться} (stydítsya)}{стыдиться (stydítsya)}]{\ru{стыдиться} (stydítsya) — to be ashamed}

\label{lesson:RU-C470-styditsya}



\begin{quote}
Before the new one: say the Russian for to envy, then the Russian for shame.
\end{quote}

\subsection*{You'll want to know: \ru{стыдиться}}
\textbf{\ru{стыдиться}} (\emph{stydítsya}) — \textquotedblleft{}to be ashamed\textquotedblright{}.

\begin{quote}
\ru{Он} \ru{стыдится} \ru{своей} \ru{ошибки}.
\end{quote}

\begin{quote}
\emph{On stydítsya svoyéy oshíbki.}
\end{quote}

\begin{quote}
He is ashamed of his mistake.
\end{quote}

Imperfective; present \ru{стыжусь}, \ru{стыдишься} (\ru{д} becomes \ru{ж}). Takes the genitive: \emph{\ru{стыдиться} \ru{ошибки}}. Perfective \ru{постыдиться}.

\begin{cousinweb}
\emph{\ru{стыдиться}} is built on \emph{\ru{стыд}} 'shame' with the reflexive \emph{-\ru{ся}}; \emph{\ru{стыд}} is inherited from Proto-Slavic \emph{styd\ru{ъ}} (EtymDB).
\end{cousinweb}

\begin{grammarlens}[title={What you've built}]
One more word for telling, explaining and describing at B1.
\end{grammarlens}

\subsection*{Your turn}
\begin{itemize}
\raggedright
  \item \emph{Say it:} \emph{stydítsya}
  \item \emph{Say it:} \emph{stydítsya}, once more
  \item \emph{Say it:} a sentence of your own with \emph{stydítsya}
  \item \emph{Recall:} say the Russian for to envy, then the Russian for shame, then say \emph{stydítsya} again
\end{itemize}

\begin{tcolorbox}[breakable,colback=teal!4,colframe=teal!35!black,title={Before you move on}]
What does \ru{стыдиться} mean? (\textquotedblleft{}To be ashamed\textquotedblright{}.) Say it once more.
\end{tcolorbox}
\section[\texorpdfstring{\ru{смущаться} (smushchátsya)}{смущаться (smushchátsya)}]{\ru{смущаться} (smushchátsya) — to get embarrassed}

\label{lesson:RU-C470-smushchatsya}



\begin{quote}
Before the new one: say the Russian for shame, then the Russian for to be ashamed.
\end{quote}

\subsection*{You'll want to know: \ru{смущаться}}
\textbf{\ru{смущаться}} (\emph{smushchátsya}) — \textquotedblleft{}to get embarrassed\textquotedblright{}.

\begin{quote}
\ru{Мальчик} \ru{смущается}, \ru{когда} \ru{его} \ru{хвалят}.
\end{quote}

\begin{quote}
\emph{Málchik smushcháyetsya, kogdá yegó khvályat.}
\end{quote}

\begin{quote}
The boy gets embarrassed when people praise him.
\end{quote}

Imperfective; perfective \ru{смутиться} (\ru{смущусь}, \ru{смутишься}). The noun is \ru{смущение} — embarrassment.

\begin{cousinweb}
\emph{\ru{смущаться}} is the reflexive of \emph{\ru{смущать}}, the imperfective of \emph{\ru{смутить}} 'to trouble, to confuse'; EtymDB lists these verbs without a source, and they are usually traced to the root of \emph{\ru{мутный}} 'cloudy'.
\end{cousinweb}

\begin{grammarlens}[title={What you've built}]
One more word for telling, explaining and describing at B1.
\end{grammarlens}

\subsection*{Your turn}
\begin{itemize}
\raggedright
  \item \emph{Say it:} \emph{smushchátsya}
  \item \emph{Say it:} \emph{smushchátsya}, once more
  \item \emph{Say it:} a sentence of your own with \emph{smushchátsya}
  \item \emph{Recall:} say the Russian for shame, then the Russian for to be ashamed, then say \emph{smushchátsya} again
\end{itemize}

\begin{tcolorbox}[breakable,colback=teal!4,colframe=teal!35!black,title={Before you move on}]
What does \ru{смущаться} mean? (\textquotedblleft{}To get embarrassed\textquotedblright{}.) Say it once more.
\end{tcolorbox}
\section[\texorpdfstring{\ru{краснеть} (krasnét)}{краснеть (krasnét)}]{\ru{краснеть} (krasnét) — to blush}

\label{lesson:RU-C470-krasnet}



\begin{quote}
Before the new one: say the Russian for to be ashamed, then the Russian for to get embarrassed.
\end{quote}

\subsection*{You'll want to know: \ru{краснеть}}
\textbf{\ru{краснеть}} (\emph{krasnét}) — \textquotedblleft{}to blush\textquotedblright{}.

\begin{quote}
\ru{Она} \ru{краснеет}, \ru{когда} \ru{говорит} \ru{с} \ru{ним}.
\end{quote}

\begin{quote}
\emph{Oná krasnéyet, kogdá govorít s nim.}
\end{quote}

\begin{quote}
She blushes when she talks to him.
\end{quote}

Imperfective; present \ru{краснею}, \ru{краснеешь}. Perfective \ru{покраснеть}. Also to turn red.

\begin{cousinweb}
\emph{\ru{краснеть}} is built on \emph{\ru{красный}} 'red', whose older sense 'beautiful' lives on in \emph{\ru{красивый}}; EtymDB lists the verb without a source, and the root is usually traced to Proto-Slavic \emph{krasa} 'beauty'.
\end{cousinweb}

\begin{grammarlens}[title={What you've built}]
One more word for telling, explaining and describing at B1.
\end{grammarlens}

\subsection*{Your turn}
\begin{itemize}
\raggedright
  \item \emph{Say it:} \emph{krasnét}
  \item \emph{Say it:} \emph{krasnét}, once more
  \item \emph{Say it:} a sentence of your own with \emph{krasnét}
  \item \emph{Recall:} say the Russian for to be ashamed, then the Russian for to get embarrassed, then say \emph{krasnét} again
\end{itemize}

\begin{tcolorbox}[breakable,colback=teal!4,colframe=teal!35!black,title={Before you move on}]
What does \ru{краснеть} mean? (\textquotedblleft{}To blush\textquotedblright{}.) Say it once more.
\end{tcolorbox}
\section[\texorpdfstring{\ru{застенчивый} (zasténchivyy)}{застенчивый (zasténchivyy)}]{\ru{застенчивый} (zasténchivyy) — shy}

\label{lesson:RU-C470-zastenchivyy}



\begin{quote}
Before the new one: say the Russian for to get embarrassed, then the Russian for to blush.
\end{quote}

\subsection*{You'll want to know: \ru{застенчивый}}
\textbf{\ru{застенчивый}} (\emph{zasténchivyy}) — \textquotedblleft{}shy\textquotedblright{}.

\begin{quote}
\ru{Застенчивый} \ru{студент} \ru{редко} \ru{задаёт} \ru{вопросы}.
\end{quote}

\begin{quote}
\emph{Zasténchivyy studént rédko zadayót voprósy.}
\end{quote}

\begin{quote}
The shy student rarely asks questions.
\end{quote}

Shy by nature. Short forms: \ru{застенчив}, \ru{застенчива}.

\begin{cousinweb}
\emph{\ru{застенчивый}} is usually traced to \emph{\ru{за} \ru{стеной}} 'behind a wall', a person who keeps out of sight; EtymDB lists the word without a source, so the link is not settled.
\end{cousinweb}

\begin{grammarlens}[title={What you've built}]
One more word for telling, explaining and describing at B1.
\end{grammarlens}

\subsection*{Your turn}
\begin{itemize}
\raggedright
  \item \emph{Say it:} \emph{zasténchivyy}
  \item \emph{Say it:} \emph{zasténchivyy}, once more
  \item \emph{Say it:} a sentence of your own with \emph{zasténchivyy}
  \item \emph{Recall:} say the Russian for to get embarrassed, then the Russian for to blush, then say \emph{zasténchivyy} again
\end{itemize}

\begin{tcolorbox}[breakable,colback=teal!4,colframe=teal!35!black,title={Before you move on}]
What does \ru{застенчивый} mean? (\textquotedblleft{}Shy\textquotedblright{}.) Say it once more.
\end{tcolorbox}
